% Intended LaTeX compiler: pdflatex
\documentclass[11pt]{article}
\usepackage[utf8]{inputenc}
\usepackage[T1]{fontenc}
\usepackage{graphicx}
\usepackage{longtable}
\usepackage{wrapfig}
\usepackage{rotating}
\usepackage[normalem]{ulem}
\usepackage{amsmath}
\usepackage{amssymb}
\usepackage{capt-of}
\usepackage{hyperref}
\author{root}
\date{\today}
\title{Extras}
\hypersetup{
 pdfauthor={root},
 pdftitle={Extras},
 pdfkeywords={},
 pdfsubject={},
 pdfcreator={Emacs 29.3 (Org mode 9.6.15)}, 
 pdflang={English}}
\begin{document}

\maketitle
\tableofcontents

\section{Notes here}
\label{sec:org6024e03}
One\footnote{Definition with two paragraphs.

Second paragraph of definition.\label{orgf90e91f}} two\footnote{anonymous \textbf{bold} note} three\footnote{Plain def.\label{org16e1dc4}} again\textsuperscript{\ref{orgf90e91f}} twice\textsuperscript{\ref{org16e1dc4}}.
\section{Misc}
\label{sec:orgac821d9}
Line one
line two \\[0pt]
line three

\textbf{raw}

\begin{itemize}
\item item with \href{https://x.org}{link} and \texttt{verb}
\begin{itemize}
\item nested [X] \texttt{code}
\end{itemize}
\end{itemize}

Text with a \hyperref[org7caf573]{radio} \label{org7caf573}Radio and use \hyperref[org7caf573]{radio}.
Macro  stays? and a COMMENT line:

Text after.
\end{document}
